\subsection{Primitive roots of unity}

Let \(n \geqslant 1\) be an integer. By the Euler indicator of \(n\) , written \(\varphi(n)\) , we understand the number of invertible elements of the ring \(\mathbf{Z}/n\mathbf{Z}\) of integers modulo n.~By the next proposition \(\varphi(n)\) is also the number of integers k prime to n and such that \(0 \leqslant k < n\) .

PROPOSITION 3. --- Let k and \(n \geqslant 1\) be two integers. Then the following assertions are equivalent :

\begin{enumerate}
\def\labelenumi{\alph{enumi})}
\item
  the residue class of k mod n is invertible in the ring Z/nZ ;
\item
  the residue class of k mod n generates the cyclic group Z/nZ ;
\item
  the integers \(k\) and \(n\) are relatively prime (I, p.~112).
\end{enumerate}

Each of the conditions a) and b) means that there exists an integer x such that \(kx \equiv 1 \pmod{n}\) , that is, there exist two integers x and y such that \(kx + ny = 1\) . This latter condition just means that k and n are relatively prime.

COROLLARY 1. --- Let G be a cyclic group of order n and let d be a divisor of n.~Then the number of elements of order d in G is equal to \(\varphi(d)\) . In particular, \(\varphi(n)\) is the number of generators of G.

Since G is isomorphic to Z/nZ, the number of generators of G is equal to \(\varphi(n)\) by Prop. 3. Let g be a generator of G; then the elements h of G such that \(h^{d}=1\) constitute the subgroup H of G generated by \(g^{n/d}\) ; this group is cyclic of order d and the elements of order d of G are the generators of H, hence their number is \(\varphi(d)\) .

COROLLARY 2.--- For every integer \(n \geqslant 1\) , we have

\[
\sum_ {d \mid n} \varphi (d) = n,\tag{1}
\]

where the integer \(d\) runs over the set of divisors \(>0\) of \(n\) \(^{1}\) .

With the notation of Cor. 1, every element of G has a finite order which is a divisor d of n and for a fixed d there are \(\varphi(d)\) such elements.

The calculation of \(\varphi(n)\) is based on the two formulae:

\begin{enumerate}
\def\labelenumi{(\arabic{enumi})}
\setcounter{enumi}{1}
\item
  \(\varphi(mn) = \varphi(m)\varphi(n)\) if \(m\) and \(n\) are relatively prime,
\item
  \(\varphi(p^a) = p^{a-1}(p - 1)\) ( \(p\) prime, \(a \geqslant 1\) ).
\end{enumerate}

The first follows at once from the fact that the rings \(\mathbf{Z} / mn\mathbf{Z}\) and \((\mathbf{Z} / m\mathbf{Z}) \times (\mathbf{Z} / n\mathbf{Z})\) are isomorphic (I, p.~112), and that \((\mathbf{A} \times \mathbf{B})^{*} = \mathbf{A}^{*} \times \mathbf{B}^{*}\) for two rings A and B. To prove (3) we note that the positive divisors of \(p^{a}\) are 1, \(p\) , \(p^{2}, \ldots, p^{a}\) ; hence the integer \(k\) has no common divisor with \(p^{a}\) other than 1 if and only if it is not divisible by \(p\) ; since there are \(p^{a-1}\) multiples of \(p\) between 0 and \(p^{a}-1\) , we find indeed (3).

The formulae (2) and (3) show at once that

\[
\varphi (n) = n \prod_ {p} (1 - 1 / p),\tag{4}
\]

where p runs over the set of prime divisors of n.

An n-th root of unity is said to be primitive if it is of order n; if there exists such a root \(\zeta\) , the group \(\mu_{n}(\mathbf{K})\) is of order n and is generated by \(\zeta\) . * For example, the primitive n-th roots of unity in C are the numbers \(e^{2\pi ik/n}\) with \(0 \leqslant k < n\) and k prime to n.~* Cor. 1 of Prop. 3 now implies the following result.

PROPOSITION 4. --- Let \(n \geqslant 1\) be an integer; suppose that there exist \(n\) \(n\) -th roots of unity in \(K\) (this is true for example if \(K\) is separably closed and \(n \cdot 1_{K} \neq 0\) ). Then the number of primitive \(n\) -th roots of unity in \(K\) is equal to \(\varphi(n)\) .

